\section{Status and roadmap}
\label{sec:status}

\subsection{What exists}

The measurement campaign is complete and its records are public: six phases,
more than 700 points, an expectations file frozen before each phase, a validity
review, a consolidated findings ledger and the nineteen-row anomaly table.

The golden C model exists and is partly validated. The queue core, the
transaction-level PCIe fabric, the host-memory model and the submission model
are landed and validated. The hardware profile, the anomaly table, the C
facade, the packetizer, the outstanding-work window, the transmit pacer, the
ingress meter, the receive processor and the go-back-N requester transport are
landed and validated. Rate control is landed behind an opt-in whose disabled
path is byte-identical to the previous slice: the notification point at the
receiver's own ingress meter, the per-queue-pair reaction point in front of the
transmit pacer, and the tail-drop egress queue that the test fabric needs
before either can be exercised. The internal arbiter of B17 is designed and not
landed.

The architecture of \S\ref{sec:arch} is a design, not a build. No
register-transfer-level code exists yet. \TODO{this is the honest status for
v0.1; if the first block lands before submission, replace this paragraph with
resource and timing numbers from the actual build rather than estimates.}

\subsection{Open clauses}

Three are load bearing and none of them is closed.

\textbf{The ingress ledger (B10).} The discards the campaign counted are far
too few to explain the goodput deficit they are supposed to cause. The meter
reproduces the equilibrium with a fitted drain rate; it does not explain it.
\anom{17} narrows the missing mechanism to a stall process with a
characteristic burst length that reverse traffic suppresses without shortening.
The first hardware build should instrument this before it does anything else,
because the answer is an occupancy trace that no vendor device will give us.

\textbf{Two anchors, one latent.} The drain rate fitted from the lone-flow
anchors makes the congestion-control loop settle below the switch egress rate,
which makes the incast tax come out near zero instead of at the measured
26.9\,\%. The two anchors want different values of the same parameter. This is
the same shape as the duplex-versus-simplex incompatibility already registered
against the same block, and we expect one mechanism to resolve both.

\textbf{The notification ledger.} Senders handle $2.24\times$ the notifications
the receiver reports sending, reproducibly, across firmware revisions and load
levels. We have no explanation, the architecture publishes both counters
without reconciling them, and we would welcome one.

\subsection{Roadmap}

\begin{enumerate}
  \item Land the internal arbiter (B17) in the golden model and close
    \anom{05}, which is the last unmodelled emergent row that has a known
    mechanism.
  \item Build B1, B2, B3, B16 and B19 first, that is, the host interface, the
    submission path and the observation path. This subset is testable against
    the golden model with no wire at all and establishes the register map and
    the monitoring ring before any transport logic depends on them.
  \item Build B9, B10 and B11 and instrument the ingress meter. Run the
    lone-flow and drain-window experiments against the implementation and
    against the real adapter side by side.
  \item Build the transport blocks (B7, B8, B12, B13) and close the anomaly
    rows that depend on them.
  \item Build the congestion-control blocks (B14, B15) last, because their
    acceptance data is the slowest to collect and depends on every block below
    them being right.
  \item Run a work-queue publication campaign to split $T_{\text{eff}}$ into
    its five stages, which converts the largest single \evcal{} parameter in
    the model into five measured ones.
  \item Repeat the campaign against a ConnectX-7 adapter at 400\,GbE to test
    the \evdecl{} scaling assumptions, every one of which is currently a
    multiplication by four.
\end{enumerate}

\section{Conclusion}
\label{sec:conclusion}

A commodity RDMA NIC can be reconstructed, from public information and
black-box measurement alone, into an architecture detailed enough to build and
honest enough to trust. The reconstruction is not a guess at the vendor's
design; it is a set of blocks, each of which names the measurement that
constrains it and the class of evidence behind that constraint, connected by a
data path that moves bytes and a monitoring network that moves nothing but
observations.

The measurements that made it possible were mostly not the ones we expected to
take. The result we are most confident of, the fixed-offset initiation law with
$T_{\text{eff}} = 4.48$\uS{} and $C = 97.1$\Gbps{}, came from an arm that was
added only after we discovered that the tool flag we thought set queue depth
did not. The result that changed the architecture most, that the congestion
notification point is the receiving adapter rather than the switch, contradicted
two of our own earlier inferences and one site belief. The result that will
matter most to a serving simulator, that inserting a gap can raise goodput by
13.8\,\%, is the opposite of what pacing is supposed to do. Freezing each
phase's expectations before running it is what let us treat all three as
findings rather than as mistakes, and it is the part of the method we would
recommend to anyone doing this again.

What remains open is stated as openly as what is closed. One block carries a
fitted stand-in for a mechanism nobody has identified, two independent
measurements want different values for the same parameter, and one counter
ledger does not balance. An implementation with its observation path designed in
from the beginning is the instrument that can settle all three, which is the
main reason to build it.
